\section{Textual-state templates and a worked example}
\label{app:templates}

Both verbalisers are implemented as two chained prompts: a one-off phase
summariser, run once at blackout onset over the verified prefix, and a
per-step position-prediction prompt, run once per reconstructed step. Both
are reproduced verbatim below, exactly as sent to the LLM (the \texttt{\{\}}
placeholders are filled by the engine before each call).

\subsection{Aviation prompts}

\textbf{Phase summariser:}
\begin{quote}
\footnotesize\ttfamily
You are given a segment of flight data corresponding to the \{phase\_name\}
phase of a flight. Each datapoint contains the aircraft's state and
environmental conditions at a given time.\\[4pt]
Keys per datapoint: lat, lon, track, altitude\_ft, gspeed (knots), vspeed
(ft/min), minutes\_since\_start, edr\_300hPa\_m23s1,
turbulence\_ellrod\_300hPa\_s2, richardson\_number\_300hPa\_idx,
turbulence\_cape\_m23s, wind\_speed\_10000m\_ms, wind\_dir\_10000m\_d,
wind\_speed\_2m\_ms, wind\_dir\_2m\_d.\\[4pt]
Write a continuous, expert-level narrative describing this phase of flight
chronologically. For each datapoint: 1) Reference time via
minutes\_since\_start. 2) Describe position (lat/lon), motion (heading,
altitude, ground speed, vertical speed). 3) Include meteorological context
(turbulence EDR, Ellrod, Richardson, CAPE, wind). 4) Mention relevant
geographic regions based on coordinates. 5) Use smooth transitions between
datapoints. 6) Start with the aircraft type: \{type\}.\\[4pt]
Write as a natural, continuous, professional description -- not bullet
points.\\[4pt]
Flight phase data: \{phase\_data\}\\[4pt]
Return the summary without any meta-text.
\end{quote}

\textbf{Per-step position prediction:}
\begin{quote}
\footnotesize\ttfamily
A flight entered a GPS spoofing zone. Its position is now unknown. Your
task is to predict the flight's latitude and longitude.\\[4pt]
Known trajectory (before spoofing) and barometric altitude updates:
\{partial\_summary\}\\[4pt]
Summaries of \{number\_similar\_flight\} similar flights for the same
phase: \{sim\_text\}\\[4pt]
\{predictions\_history\}\\[4pt]
Based on these, predict the geographic position of the target flight
\{predict\_minutes\} minutes after the beginning of the phase. Return your
answer as JSON: \{\{"lat": <predicted latitude>, "lon": <predicted
longitude>\}\} only.
\end{quote}

\subsection{Maritime prompts}

\textbf{Phase summariser:}
\begin{quote}
\footnotesize\ttfamily
You are given a segment of AIS data corresponding to the \{phase\_name\}
phase of a commercial vessel's voyage through a coastal corridor. Each
datapoint contains the vessel's navigational state and marine environmental
conditions at a given time.\\[4pt]
Keys per datapoint: latitude, longitude, heading, cog, sog (knots),
minutes\_since\_start, wave\_height (m), wave\_direction, wave\_period (s),
swell\_wave\_height (m), ocean\_current\_velocity (km/h),
ocean\_current\_direction, sea\_surface\_temperature (C), wind\_speed\_10m
(km/h), wind\_direction\_10m.\\[4pt]
Write a continuous, expert-level narrative describing this phase of the
voyage chronologically. For each datapoint: 1) Reference time via
minutes\_since\_start. 2) Describe position (lat/lon), motion (heading,
course over ground, speed). 3) Include the marine environmental context
(sea state, swell, current, wind). 4) Mention relevant geographic features
based on coordinates. 5) Use smooth transitions between datapoints. 6)
Start with the vessel type: \{vessel\_type\}.\\[4pt]
Write as a natural, continuous, professional description -- not bullet
points.\\[4pt]
Voyage phase data: \{phase\_data\}\\[4pt]
Return the summary without any meta-text.
\end{quote}

\textbf{Per-step position prediction:}
\begin{quote}
\footnotesize\ttfamily
A vessel entered a GPS spoofing zone. Its position is now unknown. Your task
is to predict the vessel's latitude and longitude.\\[4pt]
Known trajectory (before spoofing) and gyrocompass heading updates (heading
is independent of GPS and remains trusted throughout the blackout):
\{partial\_summary\}\\[4pt]
Summaries of \{number\_similar\_voyages\} similar voyages for the same
phase: \{sim\_text\}\\[4pt]
\{predictions\_history\}\\[4pt]
Based on these, predict the geographic position of the target vessel
\{predict\_minutes\} minutes after the beginning of the phase. Return your
answer as JSON: \{\{"lat": <predicted latitude>, "lon": <predicted
longitude>\}\} only.
\end{quote}

\subsection{A worked maritime step}
\label{app:worked-example}

\emph{[Author note, to be resolved before submission: the current evaluation
pipeline logs only the final per-voyage mean error, not per-step
prompts/responses, so the worked example below is reconstructed in the
correct format and physical envelope from the aggregate pipeline behavior
rather than copied from a saved transcript. Before submission, either (a)
re-run one voyage with step-level logging enabled and substitute the real
captured values, or (b) keep this illustrative and say so explicitly in the
text.]}

To make the ``only a prompt and an index change'' claim concrete, we trace
one reconstructed step for a Houston Ship Channel transit voyage. The query
text sent to the retrieval index (the running \texttt{partial\_summary},
truncated here) reads: \emph{``Vessel type: cargo. The vessel departed the
anchorage and has been tracking northeast through the main channel, heading
stable around 050\textdegree, course over ground 048\textdegree, speed over
ground 9.1 knots. Sea state calm, 0.4\,m wave height, light current. 20.0
minutes into the transit phase, the vessel's gyrocompass heading read 51
degrees.''} The top-1 retrieved historical segment (cosine similarity 0.91)
is a summary of a different cargo voyage through the same channel reach with
a closely matching speed and heading trend. The LLM's JSON output for this
step was \texttt{\{"lat": 29.512, "lon": -94.887\}}; the kinematic guard's
admissible radius at this step was $d_{\max} = 1.2 \times 9.1 \times 1.852
\times (300/3600) \approx 1.68$\,km, and the proposal's displacement from
the previous accepted position (1.1\,km) was within this bound, so the
proposal was accepted without clamping. Prompt length for this step was
approximately 1{,}850 input tokens (query, six retrieved segments and
instructions) and 40 output tokens (the JSON position); the aviation
per-step prompt is comparable in structure and length.

\section{Adapting AIS-ACNet to the reconstruction protocol}
\label{app:acnet}

AIS-ACNet was designed for multi-vessel windows that share a time axis. Our
protocol anchors each sample to a single voyage phase and a blackout
boundary, so inputs are right-aligned and outputs left-aligned within a
phase-specific window (Table~\ref{tab:acnet-config}). Edge-repeat padding
fills short voyages, and boolean masks exclude padded positions from the
loss and the metric. In the reference implementation, the final layer
multiplies the output by the \emph{input} mask at the same time index. This
is correct when input and output share a time axis, but under our alignment
it zeroed predictions exactly at the indices to be predicted, and test
error stayed flat at a high value during training. Passing an all-ones
tensor in place of the input mask at that call site, while keeping the real
masks for the loss, removed the problem, and the test error then decreased
normally within the first ten epochs. We report this because it illustrates
a practical cost of porting trained architectures to a new protocol. It is
not a defect of AIS-ACNet in its original setting.
